% ============================================================
% DOE Genesis Mission CFP 2026 — Pre-Proposal
% PA-AKG: Provenance-Aware Agentic Knowledge Graphs
% Focus Area: 19B — Hypothesis Generation from Multi-Modal Data
% ============================================================

\documentclass{files/proposalnsf}

\usepackage[utf8]{inputenc}
\usepackage[T1]{fontenc}
\usepackage{graphicx}
\usepackage{hyperref}
\usepackage{microtype}
\usepackage{enumitem}
\usepackage{wrapfig}
\usepackage[footnotesize]{caption}
\setlist{nosep, leftmargin=1.5em, topsep=0.2em}

\usepackage{xspace}
\usepackage{xcolor}

\newcommand{\todo}[1]{{\color{red}\bfseries [[TODO: #1]]}}

\newcommand{\tool}{\textcolor{blue}{\texttt{cool\_tool}}\xspace}

\newcommand{\approach}{\textcolor{blue}{\textit{approach}}\xspace}

\newcommand{\proposalEleven}{\textbf{Provenance-Aware Agentic Knowledge Graphs}}
\newcommand{\proposalEighteen}{\textbf{Provenance-Aware Agentic Knowledge Graphs}}
\newcommand{\proposalNineteen}{\textbf{Provenance-Aware Agentic Knowledge Graphs}}


\newcommand{\ie}{\textit{i.e.,}\xspace}

\newcommand{\eg}{\textit{e.g.,}\xspace}

\newcommand{\etal}{\textit{et al.}\xspace}

\begin{document}

\begin{center}
    \large{\proposalNineteen} \\ \vspace{0.5pt}
\footnotesize{PI: Chris Brown, Computer Science, Virginia Tech} \\
\footnotesize{Co-PI: Xuan Wang, Computer Science, Virginia Tech} \\
\footnotesize{Co-PI: Scott T. Steinmetz, Sandia National Labs} \\
\footnotesize{GRA: Jason Cusati, Computer Science, Virginia Tech}
\end{center}

\small{\noindent\textbf{Focus Area:} 19 -- AI for Scientific Reasoning} \hfill \\
\small{\textbf{Topic:} B. Hypothesis Generation from Multi-Modal Data}

\subsection*{Project Summary / Abstract}

The pace of scientific discovery is rapidly increasing, fueled by growing access to heterogeneous data sources---scientific literature, experimental datasets, simulation outputs, and visual data---across DOE research domains. Recent advances in large language models (LLMs) and generative AI have produced systems capable of synthesizing vast bodies of scientific content to support research. Yet AI systems cannot systematically identify gaps in scientific knowledge or generate novel, testable hypotheses with traceable provenance connecting back to the specific multi-modal evidence that supports each inference. Without persistent knowledge infrastructure, every agent session starts from scratch---unable to build on prior findings, detect genuine gaps across accumulated evidence, or produce hypotheses that human researchers can audit, extend, or act on---resulting in a fragmented AI-driven discovery cycle.

To address these challenges, we will introduce \textbf{provenance-aware agentic knowledge graphs} (PA-AKG) as foundational infrastructure for AI-driven scientific hypothesis generation from multi-modal data. A persistent research knowledge graph (KG) encodes the current state of scientific understanding---problems, datasets, methods, constraints, and evidence---linked by explicit semantic relations, and is continuously updated as new evidence emerges from literature, experiments, simulation outputs, and visual data. LLM-based agents compare incoming findings against KG-encoded knowledge, identify under-evidenced gaps, and generate structured, evidence-grounded hypotheses carrying explicit provenance chains across all ingested modalities. PA-AKG will offer a significant AI advantage over existing retrieval-augmented approaches, enabling systematic gap identification and traceable hypothesis generation at a scale and fidelity not achievable by human researchers alone.

The proposed work will design and implement the PA-AKG infrastructure, incorporating AI directly into multi-modal hypothesis generation workflows for data collection, integration, analysis, and curation. We will evaluate PA-AKG against existing approaches, investigating the traceability, scientific quality, and generalizability of AI-generated hypotheses. This project will advance the foundations for AI-driven scientific reasoning by designing techniques to persist and integrate knowledge across heterogeneous data modalities, enabling AI to reason over an evolving, multi-source model of scientific understanding. This work will produce deployable tools, including an open-source PA-AKG implementation, multi-modal evaluation benchmarks, and empirical insights to lay the groundwork for a new generation of trustworthy AI-augmented scientific hypothesis generation.

\newpage
\pagenumbering{arabic}

\begin{center}
    \large{\proposalEleven} \\ \vspace{0.5pt}
\end{center}\vspace{-15pt}

\subsection*{Cover Sheet} \vspace{12pt}

\textbf{Lead Applicant/Institution:} Virginia Tech \\

\noindent\textbf{Street Address:} 925 Prices Fork Road, Blacksburg, VA 24060 \\

\noindent\textbf{Postal Address:} 150 Kent St., Blacksburg, VA 24061\\

\noindent\textbf{Lead PI:} Dr. Chris Brown, 540-231-9176, dcbrown@vt.edu \\

\noindent\textbf{Administrative Point of Contact:} Robert Marcum, (540) 231-7537, marcumr@vt.edu \\

\noindent\textbf{RFA Number:} DE-FOA-0003612 \\

\noindent\textbf{Primary Challenge:} \\

\noindent
 19 -- AI for Scientific Reasoning \\
B. Hypothesis Generation from Multi-Modal Data \\

\begin{table}[h]
    \begin{tabular}{|l|l|}
    \hline
    \multicolumn{2}{|c|}{\textbf{Senior/Key Personnel for Lead Institution and All Partner Institutions}}  \\ \hline
        \textbf{Institution Name} & \textbf{Senior/Key Personnel Including the Director} \\ \hline
        Virginia Tech & PI Brown; Co-PI Wang \\ \hline
        Sandia National Labs & Co-PI Steinmetz \\ \hline
    \end{tabular}
    \label{tab:my_label}
\end{table}

\begin{table}[h]
    \begin{tabular}{|l|l|}
    \hline
    \multicolumn{2}{|c|}{\textbf{Summary budget information for lead institution and all partner institutions}}    \\ \hline
    \textbf{Institution} & \textbf{Year 1 Budget \$} \\ \hline
        Virginia Tech (lead) & \todo{} \\ \hline
        Sandia National Labs &  \todo{} \\ \hline
    \end{tabular}
    \label{tab:my_label}
\end{table}

\newpage
\pagenumbering{arabic}

\begin{center}
    \large{\proposalNineteen} \\ \vspace{0.5pt}
\footnotesize{PI: Chris Brown, Computer Science, Virginia Tech} \\
\footnotesize{Co-PI: Xuan Wang, Computer Science, Virginia Tech} \\
\footnotesize{Co-PI: Scott T. Steinmetz, Sandia National Labs} \\
\footnotesize{GRA: Jason Cusati, Computer Science, Virginia Tech}
\end{center}\vspace{-15pt}

\subsection*{Project Narrative}

\paragraph*{Introduction}
The volume of scientific data and heterogeneous knowledge sources are expanding at an unprecedented rate~\cite{perabytes,steiner2023exponential}. For instance, experimental results from facilities such as the Advanced Photon Source~\cite{galayda1995advanced}, simulation outputs from exascale computing platforms, and domain literature---spanning materials science, energy research, and national security---are accumulating faster than researchers can manually synthesize and act on them. This widening \textbf{knowledge synthesis gap} poses significant challenges for scientific discovery: identifying genuine knowledge gaps, detecting cross-modal correlations, and formulating novel testable hypotheses requires integrating evidence across sources that no individual researcher can read or process simultaneously~\cite{zenodo1183562}, leading to duplication of effort and wasted resources.

Recent advances in large language models (LLMs) and generative AI have produced tools of remarkable capability to support scientific reasoning. LLM-based systems can retrieve and summarize vast bodies of literature~\cite{hsu2024chime} and generate plausible scientific text~\cite{huang2023role} and hypotheses~\cite{bazgir2025agentichypothesis} across domains and input types (\eg natural language chats, documents, images, etc.). The potential of LLM-based agents to accelerate DOE-relevant discovery is widely recognized and actively being pursued across the national laboratory system~\cite{rummler2026sandians}. However, LLMs are prone to hallucinate---generating plausible-sounding scientific content without the ability to distinguish between well-supported claims and speculative or incorrect ones~\cite{Ji2023HallucinationSurvey}---including fabricated citations, incorrectly recalled experimental results, and duplicate hypothesis generation. These outcomes can misdirect research programs and waste substantial resources. Recently, agentic systems coordinating multiple AI agents~\cite{saeedi2025astroagents} and augmentation methods to retrieve knowledge from relevant documents~\cite{} aim to generate hypotheses and advance scientific discovery through grounding and contextualizing model reasoning. However, fundamental challenges that constrain LLM utility for rigorous, reproducible hypothesis generation remain:

\textbf{Absence of persistent knowledge state.} Current LLM-based systems lack the ability to recall prior sessions or the accumulated state of a research program. This prevents AI from building cumulative understanding across heterogeneous sources, detecting contradictions against established prior findings, and avoiding redundant hypothesis generation. As a result, agents provide informative but not progressive insights---failing to build on prior evidence and re-deriving the same correlations session after session.

\textbf{Lack of traceable provenance.} Scientific hypotheses must be auditable~\cite{}. When an AI system generates a hypothesis from multi-modal evidence---combining findings from literature, experimental data, and simulation outputs---researchers must be able to trace that hypothesis to the specific sources that support each inferential step. While current systems may surface relevant documents, they do not provide structured reasoning that connects multi-modal evidence to conclusions. Without explicit provenance, AI-generated hypotheses cannot be rigorously evaluated, extended, or reliably used to guide further experimentation. \vspace{-5pt}


\problem{The lack of structured and persistent multi-modal knowledge infrastructure limits AI agents' ability to systematically identify knowledge gaps and generate novel, traceable hypotheses for DOE scientific discovery.}{}

To address these challenges, this project proposes the design, implementation, and evaluation of \textbf{provenance-aware agentic knowledge graphs} (PA-AKG) for scientific hypothesis generation from multi-modal data. PA-AKG is foundational infrastructure for closed-loop, AI-integrated scientific discovery, providing agents with a persistent, structured, and auditable representation of the current state of scientific understanding in a given domain. Knowledge graphs (KGs) have been applied to research as a mechanism for representing complex, heterogeneous information through explicit semantic relationships~\cite{bai2024advancing,d2025role,karras2024organizing}. PA-AKG will encode scientific knowledge (\ie problems, datasets, experimental methods, theoretical frameworks, and empirical observations from literature, simulation, and visual sources) as nodes, with relationships between them as directed edges. As new evidence emerges across modalities, LLM-based agents compare incoming findings against KG-encoded knowledge, identify under-evidenced problem nodes, and generate structured, evidence-grounded hypotheses---each carrying an explicit provenance chain that traces back to specific multi-modal sources of motivation. Confirmed hypotheses are written back to the KG, updating the current model of scientific understanding and enabling the next iteration of the discovery cycle.

\begin{wrapfigure}[15]{r}{0.41\textwidth}
    \centering
    \includegraphics[width=0.39\textwidth]{abstract_19B/images/architecture-diagram.png}
    \caption{PA-AKG Architecture Overview}
    \label{fig:pa-akg}
\end{wrapfigure}

\paragraph{Objectives.} Our goal is to design and implement a provenance-first agentic KG architecture with three integrated layers (see Figure~\ref{fig:pa-akg}). We will implement the PA-AKG architecture, developing infrastructure to ingest and structure claim-evidence pairs from heterogeneous scientific sources---text, experimental data, simulation outputs, and visual data---with role-specific agents. The initial implementation will focus on the software engineering research domain, following our prior work~\cite{cusati2026fose}, within the context of HPC. Software engineering outlines processes and methods for developing and maintaining software systems, and is also critical for technological innovation~\cite{vos2026reclaiming} and enabling the development and maintenance of applications~\cite{schmidberger2012need}. We will conduct a comprehensive evaluation of our approach, conducting quantitative benchmarking studies against baseline approaches and qualitative user studies to investigate provenance completeness, traceability, hypothesis quality, and research perceptions. 

\paragraph{Research Plan.} Our research plan involves four main phases:

\textit{Phase 1 (Months 1--3): Multi-Modal Knowledge Representation.}
We will design the research KG schema with entity types such as \textsc{Problem}, \textsc{Dataset}, \textsc{Method}, \textsc{Constraint}, \textsc{Evidence}, and \textsc{Hypothesis}---linked by semantic relations (\textsc{supports}, \textsc{contradicts}, \textsc{enables}, \textsc{requires}). We will conduct pilot studies (\eg surveys and interviews) to obtain insights from researchers across domains, informing entity types and relationships for our initial implementation. Research problems will be modeled as tractable, scoped questions with identifiable assumptions, methods, and evaluation criteria, extracted from source documents with evidence spans anchored to passages or image regions. A hybrid symbolic--semantic design pairs a property graph store with a vector index. We will populate an initial KG from the software engineering research corpus established in prior work~\cite{cusati2026fose, brown2025esem}, validate entity extraction accuracy against manually annotated ground-truth annotations, and establish baseline metrics for knowledge gap coverage and cross-modal entity resolution.

\textit{Phase 2 (Months 4--6): Multi-Modal Extraction Pipeline and Agentic Orchestration.}
We will implement the extraction and automation pipeline: structured prompting will extract claim--evidence pairs from text, data tables, and image captions; span-level alignment will link each claim to character-level text spans or region-level image references; and canonicalization auditing will normalize, validate, and deduplicate entity references across modalities to canonical KG nodes. The agentic orchestration layer, adapting the Planning \& Control pattern~\cite{Park2023GenerativeAgents}, introduces four role-specific agent types. \textit{Ranking agents} score candidate research directions by tractability and evidence strength across all ingested modalities. \textit{Continuation agents} propose follow-on analyses grounded in KG gaps, targeting under-evidenced \textsc{Problem} nodes across source types. \textit{Evaluation agents} execute reproducible analytical workflows and write results back to the KG as structured evidence nodes. \textit{Synthesis agents} generate hypothesis artifacts---problem statement, multi-modal supporting evidence chain, proposed test, predicted outcome---and curate them to the KG as \textsc{Hypothesis} nodes with full provenance metadata. The provenance inheritance constraint is enforced at every agent handoff.

\textit{Phase 3 (Months 7--9): Evaluation and Generalization.}
We will conduct two comprehensive evaluations of our PA-AKG approach. The first will benchmark our system against existing approaches, including vector-based RAG, hybrid retrieval, and GraphRAG~\cite{Edge2024GraphRAG} reasoning. We will construct a benchmark of human-annotated claim--evidence pairs from research corpora. We will calculate span-level alignment precision/recall, Cohen's $\kappa$ with expert-labeled canonical entity clusters, and traceability completeness (\ie \% of KG nodes with complete provenance chains). A secondary scientific domain will be evaluated to assess cross-domain generalizability of the multi-modal extraction and hypothesis generation pipeline. Our second evaluation will assess the quality of PA-AKG-generated hypotheses---rated by domain experts on specificity, falsifiability, and cross-modal grounding---compared against retrieval-only and graph-only baselines.

\textit{Phase 4: Dissemination.}
We will pursue targeted dissemination of PA-AKG research findings to academic and practitioner audiences. Academic outputs will be submitted to premier venues in software engineering, AI, and machine learning, including the International Conference on Software Engineering (ICSE), the Conference on Neural Information Processing Systems (NeurIPS), International Conference on Machine Learning (ICML), and the AAAI Conference on Artificial Intelligence. We will submit to tool, demo, and tutorial tracks at these venues to provide hands-on exposure of PA-AKG and its multi-modal hypothesis generation capabilities to active research communities, lowering the barrier to adoption across DOE and academic research programs. To reach the practitioner community, we will release the PA-AKG implementation as an open-source package with comprehensive documentation, worked examples, and integration guides targeting scientific researchers across DOE domains. We will engage directly with Sandia National Laboratories and partner institutions to present findings, gather feedback on multi-modal canonicalization and hypothesis quality from domain experts, and assess adoption pathways within active DOE scientific workflows. This phase will conclude with an impact assessment evaluating how PA-AKG changes researchers' ability to identify genuine knowledge gaps and generate traceable hypotheses from heterogeneous evidence---informing the design of future AI-driven scientific reasoning systems and establishing PA-AKG as a reference model for trustworthy, multi-modal AI hypothesis generation.

\paragraph{Expected Outcomes.} The expected outcomes include:
\begin{itemize}
    \item A design for provenance-aware knowledge graph architecture for AI-driven hypothesis generation from multi-modal scientific data.
    \item An open-source PA-AKG implementation, consisting of a multi-modal extraction pipeline and agentic orchestration, to persist and integrate knowledge across heterogeneous scientific sources.
    \item Multi-modal benchmarking infrastructure to assess LLM-based systems for scientific hypothesis generation.
    \item Empirical insights on the traceability, scientific quality, and generalizability of AI-generated hypotheses from multi-modal evidence.
    \item Broader scientific impact through open engagement with industry and academic research communities, informing the design of next-generation AI-driven scientific reasoning systems and establishing PA-AKG as a reference model for trustworthy, multi-modal hypothesis generation in DOE's scientific mission.
\end{itemize}

\paragraph{Research Team.} This project aligns with the prior work and trajectory of the research team.
\textbf{PI Brown} and GRA Cusati explored how LLMs propagate research-based evidence~\cite{brown2025esem} and outlined researchers' desire for structured knowledge accumulation~\cite{cusati2026fose} in software engineering contexts. Dr. Brown will lead the overall execution of the project---including the PA-AKG implementation and research plan. \textbf{Co-PI Wang} investigates LLMs in science and society, with substantial experience studying LLMs (\ie~\cite{bi2024ai,srivastava2025llmthinkbench}) and agents (\ie~\cite{lu2024triageagent,latimer2025hindsight}) in scientific settings and earning an NSF CAREER award to extract patterns from research data~\cite{wang_career}. Dr. Wang will inform the design of the PA-AKG architecture, particularly focusing on multi-modal knowledge retrieval and extraction, multi-agent orchestration, and knowledge persistence, and the benchmark evaluation~\cite{srivastava2025beyondbench}. \textbf{Co-PI Steinmetz} is a researcher at Sandia National Labs with expertise in human-machine teaming~\cite{chance2023towards,powell2021coordination}, generative AI~\cite{krofcheck2023overview,layton2025mstd}, and trust in AI systems~\cite{steinmetz2025trust,naugle2025trusted}. Dr. Steinmetz will inform the PA-AKG design and support our research efforts to evaluate the usability and trustworthiness of our approach against DOE scientific literature and experimental data corpora in practical research contexts. \textbf{GRA Cusati} is a Ph.D. student in the Department of Computer Science at Virginia Tech with extensive industry experience implementing complex production software in various roles. Jason will lead the design and development of the PA-AKG infrastructure and the overall data collection and analysis efforts.

\paragraph{Resources.} Virginia Tech's Advanced Research Computing provides GPU cluster infrastructure for large-scale KG construction and LLM inference workloads. The team will access publicly available scientific literature corpora (\eg ACM Digital Library, arXiv, IEEE Xplore) and pre-trained language models (\eg Llama, GPT-4, etc.) for multi-modal extraction and orchestration. Sandia National Laboratories will contribute access to DOE-relevant scientific literature and experimental data corpora for secondary domain validation, ensuring the research addresses realistic national laboratory knowledge synthesis challenges.


\bibliography{references}
\bibliographystyle{abbrv}

\end{document}
